\chapter*{Kivonat}
\paragraph{}A dolgozat azt vizsgálja, hogy a kereskedelmi API-n elérhető multimodális nagy nyelvi 
modellek mennyire pontosan becsülik egy étkezés tömegét, energia- és makrotápanyag-tartalmát 
fénykép, szöveges leírás vagy a kettő együttes felhasználásával, egy cukorbetegség kezelését támogató alkalmazás 
étkezésnaplózó funkciójához. Öt képes és egy szöveges nyilvános adathalmazon egységes mérési 
keretrendszerrel hasonlítottam össze a Google, az OpenAI és az Anthropic modelljeit több árszinten, 
valamint promptváltozatokat, több nézőpontból készült képeket és szöveges bemenetet. Az 50 ételes 
előmérések alapján előre rögzített döntési mérést végeztem, amelynek hipotéziseit az előmérésben 
fel nem használt ételeken, Holm-korrekcióval vizsgáltam. A pontosság mellett a költséget, a 
válaszidőt, a strukturált kimenet megbízhatóságát, a válaszok konzisztenciáját és az ételek 
megnevezését is értékeltem.

\paragraph{}A döntési mérés modelljei közül a legolcsóbb és leggyorsabb, a Gemini Flash-Lite szénhidrátban a 
Nutrition5k, az OmniFood8k és az ACETADA adathalmazon szignifikánsan pontosabb volt a nagyjából 
öt-hatszor drágább Claude Sonnet és GPT-6.1 Sol modelleknél, a DFood-Tj-n viszont a drágább modellek 
voltak pontosabbak. A promptváltozatok szénhidrátban nem javítottak következetesen, az étel szokásos 
elkészítési módját is figyelembe vevő P3\_recipe prompt viszont az OmniFood8k-n csökkentette a 
zsírbecslés hibáját. A döntési mérés 5546 hívása között nem volt sikertelen hívás, a sémának 
megfelelő válaszok között azonban nagyságrendi hibák is előfordultak, és a becslés hibája 
ételenként nagy lehet. Ezért a becslést egy alkalmazásban a felhasználó által ellenőrizhető és 
javítható javaslatként érdemes megjeleníteni.
\clearpage
